\hypertarget{new-unit-request-template-request-template-draft-a-new-hackstreak-pack}{%
\subsection{Request template: draft a new hackstreak
pack}\label{new-unit-request-template-request-template-draft-a-new-hackstreak-pack}}

Copy everything below the line into a NEW chat with Claude (it does not
remember earlier chats, so this carries everything). Fill in the five
fields marked \textless\textless...\textgreater\textgreater. Keep the
rest as it is.

\begin{center}\rule{0.5\linewidth}{0.5pt}\end{center}

Draft a hackstreak pack (a learning roadmap) for me. Reply with ONLY the
TOML file, nothing else.

\hypertarget{new-unit-request-template-my-request}{%
\subsubsection{My request}\label{new-unit-request-template-my-request}}

\begin{itemize}
\tightlist
\item
  Topic / unit: \textless\textless e.g. "Networking: how packets really
  move, with Beej\textquotesingle s guide and
  Wireshark"\textgreater\textgreater{}
\item
  Kind: \textless\textless"proficiency" (I want to get genuinely good at
  this subject, give me 8-15 levels) OR "project" (one concrete build
  with a finish line, 3-6 levels) -\/- pick
  one\textgreater\textgreater{}
\item
  What I already know or have finished: \textless\textless e.g. "Bandit
  0-12, basic C, Python, ssh"\textgreater\textgreater{}
\item
  Time reality: \textless\textless e.g. "10-30 min on deadline days, 1-2
  h on lab nights, 4+ h on free half-days"\textgreater\textgreater{}
\item
  Hardware / tools: \textless\textless e.g. "Ubuntu 24.04 + Windows 11
  dual boot, RTX 4060 8 GB, 16 GB RAM, Raspberry Pi 5 8 GB, LAN
  cable"\textgreater\textgreater{}
\item
  What real proficiency/finished should look like:
  \textless\textless e.g. "I can capture and explain a TCP handshake,
  write a raw-socket tool, and debug a real connectivity issue without
  googling the basics"\textgreater\textgreater{}
\end{itemize}

\hypertarget{new-unit-request-template-about-me-stable-keep-as-is}{%
\subsubsection{About me (stable, keep as
is)}\label{new-unit-request-template-about-me-stable-keep-as-is}}

Wide interests (software dev, web, mobile, game dev, OS, networks,
low-level, cybersecurity) and I get overwhelmed and lose the thread when
a plan is open-ended. I distrust AI-generated learning content. I only
want real, free, well-known resources, and I learn best from small
puzzles with a real check. I am a CSE student at an IIT with heavy
academic load. I only attack machines I own.

\hypertarget{new-unit-request-template-format-rules-the-hackstreak-pack-format}{%
\subsubsection{Format rules (the hackstreak pack
format)}\label{new-unit-request-template-format-rules-the-hackstreak-pack-format}}

\begin{itemize}
\tightlist
\item
  Sections: \texttt{{[}pack{]}} (id, title, kind = "proficiency" or
  "project", version = 1, source = "ai-draft", reviewed = false),
  \texttt{{[}vars{]}}, \texttt{{[}{[}level{]}{]}},
  \texttt{{[}{[}node{]}{]}}.
\item
  Levels, \texttt{kind = "proficiency"}: 8-15 on the \texttt{main}
  lane, each a real sub-skill with 4-8 nodes (like a Duolingo unit), and
  each level ends with one node marked \texttt{milestone = true} (size
  M or L) that is a harder, integrative task -\/- something built or
  explained from scratch, not a repeat of the level\textquotesingle s
  small steps. That milestone is the actual proficiency check; do not
  skip it.
\item
  Levels, \texttt{kind = "project"}: 3-6 on the \texttt{main} lane
  total, same milestone rule, smaller overall.
\item
  Either kind: optionally ONE \texttt{side} lane level of S-sized quick
  wins for busy days. Each \texttt{{[}{[}level{]}{]}} has id, title,
  lane ("main" or "side").
\item
  Each \texttt{{[}{[}node{]}{]}} has: level, id (kebab-case, unique,
  never changes), title, size, needs (list of node ids, may be empty),
  done\_when, resource, optionally \texttt{milestone = true}, and
  optionally verify.
\item
  size: "S" = up to 30 min, "M" = 1-2 h, "L" = 3 h or more. Mix sizes so
  busy days have an S option.
\item
  done\_when: ONE concrete, checkable sentence. Prefer an artifact on my
  disk (a file, a scan, a script, a short writeup). Never "understand X"
  or "read chapter Y" alone; pair reading with something I produce.
\item
  resource: a real, well-known URL. If you are not sure a URL exists,
  write \texttt{resource = "TODO"} instead of guessing. Do NOT write
  lesson content or explanations, only the plan.
\item
  verify (optional): \texttt{\{ type = "manual" \}} (default),
  \path|{ type = "file", path = "...", min_bytes = N }|,
  or \path|{ type = "command", cmd = "..." }| (exit code 0 =
  pass; harmless read-only commands only, no sudo).
\item
  Use \path|series = { from = A, to = B }| for numbered
  material, with \texttt{\{n\}} and \texttt{\{next\}} in
  title/done\_when.
\item
  Use \texttt{{[}vars{]}} for values I will edit (IPs, usernames,
  folders) and refer to them as \texttt{\{name\}}.
\item
  Never ask me to store passwords, flags, or tokens.
\item
  The first node of the first level must be startable tonight in 30
  minutes or less.
\end{itemize}

\hypertarget{new-unit-request-template-tiny-example-of-the-style}{%
\subsubsection{Tiny example of the
style}\label{new-unit-request-template-tiny-example-of-the-style}}

\begin{verbatim}
[pack]
id = "example"
title = "Example"
version = 1
source = "ai-draft"
reviewed = false

[[level]]
id = "basics"
title = "Basics"
lane = "main"

[[node]]
level = "basics"
id = "first-capture"
title = "Capture my first packets"
size = "S"
needs = []
done_when = "a 30-second capture saved at ~/hack/pcaps/first.pcapng that contains at least one DNS query"
resource = "https://www.wireshark.org/docs/"
verify = { type = "file", path = "~/hack/pcaps/first.pcapng", min_bytes = 1000 }
\end{verbatim}

\begin{center}\rule{0.5\linewidth}{0.5pt}\end{center}

After you get the reply: save it as
\texttt{packs/\textless{}id\textgreater{}.toml}, then run
\path|ascend pack lint packs/<id>.toml|,
\path|ascend pack check-links packs/<id>.toml|,
read every \texttt{done\_when}, then
\texttt{ascend review \textless{}id\textgreater{}} and
\path|ascend pack install <id>|. (The CLI
binary is \texttt{ascend}, not \texttt{hack} --- the earlier drafts of
this project used \texttt{hack} before the public rename; use
\texttt{ascend} for every command.)
